% DOCUMENT LAYOUT
\usepackage{fontspec}
\usepackage{geometry}
\geometry{a4paper, left=3.5cm, right=1.5cm, textheight=26cm, marginparsep=6pt, marginparwidth=2.0cm, hmarginratio=3:2}
\setlength\parindent{0in}
\usepackage{wrapfig}
\usepackage{caption}
% \usepackage{siunitx}

% FONTS
\usepackage{libertine}
% \usepackage[usenames,dvipsnames]{color}
\usepackage[svgnames]{xcolor}
% \usepackage{xunicode}
\usepackage{xltxtra}
\defaultfontfeatures{Ligatures=TeX}
% \setmainfont [Ligatures={Common, Rare, Discretionary}, Numbers={OldStyle}]{Linux Libertine O}
\setmainfont [Ligatures={Common, Rare, Discretionary}, Numbers={Proportional}]{Linux Libertine O}
%\setmonofont[Scale=0.75]{Linux Libertine Mono O}
\setmonofont[Scale=0.75]{DejaVu Sans Mono}
\usepackage{fontawesome5}
\usepackage{academicons}
\usepackage{cancel}
% \usepackage{titlesec}

% ---- CUSTOM COMMANDS
% \chardef\&="E050
\newcommand{\html}[1]{\href{#1}{\scriptsize\textsc{[html]}}}
\newcommand{\pdf}[1]{\href{#1}{\scriptsize\textsc{[pdf]}}}
\newcommand{\doi}[1]{\href{#1}{\scriptsize\textsc{[doi]}}}
% ---- CUSTOM AMPERSAND
\newcommand{\amper}{{\fontspec[Scale=.95]{Gentium Basic Italic}\selectfont\itshape\&}}
% ---- LÍNEA TEMPORAL (años en el margen izquierdo, unidos por una línea vertical)
% Uso:  \begin{timeline} \years{2024} texto ... \years{2024} otro ... \end{timeline}
%        \begin{ntimeline} ... \end{ntimeline}   (igual, con numeración inversa)
% \years{etiqueta} abre un ítem; ítems consecutivos con la misma etiqueta se
% agrupan (la etiqueta y el punto grande aparecen sólo en el primero).
\usepackage{framed}
\makeatletter
\newlength\tlrulesep  \setlength\tlrulesep{9pt}   % distancia texto -- línea
\newlength\tlyearsep  \setlength\tlyearsep{7pt}   % distancia línea -- año
\newlength\tlgroupsep \setlength\tlgroupsep{3pt}  % espacio extra entre grupos
\colorlet{tlline}{DimGrey!50}
\colorlet{tldot}{MidnightBlue}
\colorlet{tlyear}{DimGrey}
\def\tl@last{}
\def\tl@subdot{\rlap{\kern-\tlrulesep\kern.25pt\makebox[0pt]{\textcolor{tlline}{\raisebox{.38ex}{\fontsize{3.5}{3.5}\selectfont\faCircle}}}}}
\def\tl@maindot#1{%
  \llap{\textcolor{tlyear}{\small #1}\kern\tlrulesep\kern\tlyearsep}%
  \rlap{\kern-\tlrulesep\kern.25pt\makebox[0pt]{\textcolor{tldot}{\raisebox{.3ex}{\fontsize{5.5}{5.5}\selectfont\faCircle}}}}}
\let\tl@marker\tl@subdot
\newcommand{\years}[1]{%
  \par
  \def\tl@this{#1}%
  \ifx\tl@this\tl@last
    \global\let\tl@marker\tl@subdot
  \else
    \ifx\tl@last\@empty\else\addvspace{\tlgroupsep}\fi
    \gdef\tl@marker{\tl@maindot{#1}}%
    \global\let\tl@last\tl@this
  \fi
  \item}
\def\tl@makelabel#1{\tl@marker\global\let\tl@marker\tl@subdot\hss\llap{#1}}
\newlength\tldoty \setlength\tldoty{3.6pt}  % altura del centro del punto sobre la línea base
\newbox\tl@boxa \newbox\tl@boxb \newbox\tl@boxc
\def\tl@rulehbox#1#2{\hbox{\hskip-\tlrulesep{\color{tlline}\vrule width .5pt #1}\hskip\tlrulesep\hskip-.5pt\box\tl@boxa}}
% Primer tramo: la línea arranca en el centro del primer punto (no en el borde superior del bloque)
\long\def\tl@firstframe#1{%
  \setbox\tl@boxa\vbox{#1}%
  \setbox\tl@boxb\copy\tl@boxa
  {\vfuzz\maxdimen \vbadness\maxdimen \global\setbox\tl@boxc\vsplit\tl@boxb to \z@}
  \tl@rulehbox{height \dimexpr\ht\tl@boxa-\ht\tl@boxc+\tldoty\relax depth \dp\tl@boxa}{}}
% Tramos intermedios/finales (continúan desde la página anterior): línea completa
\long\def\tl@midframe#1{%
  \setbox\tl@boxa\vbox{#1}%
  \tl@rulehbox{}{}}
\def\tl@frame{%
  \let\FrameCommand\tl@firstframe
  \let\FirstFrameCommand\tl@firstframe
  \let\MidFrameCommand\tl@midframe
  \let\LastFrameCommand\tl@midframe
  \OuterFrameSep=2pt
  \MakeFramed{\FrameRestore}%
  \global\let\tl@last\@empty
  \global\let\tl@marker\tl@subdot}
\newenvironment{timeline}{%
  \tl@frame
  \list{}{\leftmargin\z@ \labelwidth\z@ \labelsep\z@ \itemindent\z@ \listparindent\z@
    \topsep\z@ \partopsep\z@ \itemsep 1pt \parsep 2pt plus 1pt
    \let\makelabel\tl@makelabel}%
}{\endlist\endMakeFramed\global\@nobreakfalse}
\newenvironment{ntimeline}{%
  \tl@frame
  \etaremune
  \let\makelabel\tl@makelabel
  \topsep\z@ \partopsep\z@ \itemsep 2pt plus 1pt \parsep 2pt plus 1pt
}{\endetaremune\endMakeFramed\global\@nobreakfalse}
\makeatother

% HEADINGS
\usepackage{sectsty}
\usepackage[normalem]{ulem}
\usepackage{etoolbox}
\usepackage{needspace}
% Un título nunca queda solo al pie de página: exige espacio para unas líneas debajo
\makeatletter
\AtBeginDocument{%
  \pretocmd{\section}{\Needspace{5\baselineskip}}{}{}%
  \pretocmd{\subsection}{\if@nobreak\else\Needspace{5\baselineskip}\fi}{}{}%
  \pretocmd{\subsubsection}{\if@nobreak\else\Needspace{3\baselineskip}\fi}{}{}}
\makeatother
\sectionfont{\mdseries\upshape\Large}
\subsectionfont{\mdseries\scshape\large}
\subsubsectionfont{\mdseries\upshape\normalsize}

% Mi setup
\providecommand{\cvlang}{spanish}
\usepackage[\cvlang]{babel}
\usepackage{csquotes}
\usepackage[babel=true]{microtype}
\usepackage{tabularx}
\usepackage{booktabs}
\usepackage{amsmath}
\usepackage{etaremune}
\usepackage{fancyhdr}
\usepackage{lastpage}
% \usepackage{natbib}
% \usepackage{bibentry}
\usepackage[sorting=none, url=false, doi=true, eprint=false, isbn=false, maxnames=10]{biblatex}
\AtEveryBibitem{\clearlist{language}}
\addbibresource{publicaciones.bib}
\addbibresource{congresos.bib}
\AtEveryBibitem{\clearfield{note}}
\AtEveryCitekey{\clearfield{note}}


\pagestyle{fancy}
\lhead{}
\chead{}
\rhead{}
\lfoot{\color{DimGrey}\small{Manuel Carlevaro -- \textit{curriculum vit\ae}}}
\cfoot{}
\ifdefstring{\cvlang}{english}
  {\rfoot{\color{DimGrey}\small Page \thepage \hspace{0.05cm} of \hspace{0.05cm} \pageref{LastPage}}}
  {\rfoot{\color{DimGrey}\small Página \thepage \hspace{0.05cm} de \pageref{LastPage}}}
\renewcommand{\headrulewidth}{0pt}
\renewcommand{\footrulewidth}{0pt}

% PDF SETUP
% ---- FILL IN HERE THE DOC TITLE AND AUTHOR
\usepackage[unicode, bookmarks, colorlinks, breaklinks,
	pdftitle={Manuel Carlevaro - Curriculum Vitae},
	pdfauthor={Manuel Carlevaro},
	pdfproducer={http://carlevaro.ar}
]{hyperref}
\hypersetup{linkcolor=blue,citecolor=blue,filecolor=black,urlcolor=MidnightBlue}
\renewcommand{\labelitemi}{\faAngleRight}

% Formato de sección, subsección y subsubsección
% \titleformat{\section}
% {\normalfont\Large\bfseries}{\thesection}{1em}{}
% \titleformat{\subsection}
% {\normalfont\large\bfseries}{\thesubsection}{1em}{}
% \titleformat{\subsubsection}
% {\normalfont\tiny\bfseries}{\thesubsubsection}{1em}{}

\author{Carlos Manuel Carlevaro}
\title{Curriculum Vitae}

% ---- SHARED BUILDING BLOCKS (editar acá una sola vez)
% Año mínimo de los artículos a listar con \art (0 = todos)
\providecommand{\cvminyear}{0}

% \art{año}{clave-bib}: ítem de artículo (dentro de ntimeline)
\newcommand{\art}[2]{%
  \ifnum#1<\cvminyear\relax\else
    \years{#1} \fullcite{#2}.%
  \fi}

% Libro (idéntico en todas las versiones, salvo el idioma)
\newcommand{\libropython}{F. Batista \ifdefstring{\cvlang}{english}{and}{y} Manuel Carlevaro.
  \textbf{Python para Ciencia y Tecnolog\'ia}, 2025.
  \ifdefstring{\cvlang}{english}{Self-published}{Autoeditado}, ISBN: 978-631-00-9794-7.
  \url{https://libropython.science/}.}

% Pie de página final
\newcommand{\cvcolophon}{\vspace{0.5cm}\vfill{}
\begin{center}
  {\scriptsize \ifdefstring{\cvlang}{english}{Last update}{Última actualización}: \today\- •\-
  \ifdefstring{\cvlang}{english}{Generated by}{Generado con}
  \href{https://en.wikipedia.org/wiki/LuaTeX}{\fontspec{TeX Gyre Pagella}Lua\LaTeX} \\
  \href{https://carlevaro.ar}{https://carlevaro.ar}}
\end{center}}
